% Clase 5 --- Líneas de campo de Faraday (§3.2).
% Tres representaciones canónicas: campo uniforme (líneas paralelas
% equiespaciadas), carga puntual positiva (radiales salientes) y dipolo
% (nacen en +Q y mueren en -Q).
\begin{center}
\begin{tikzpicture}[scale=1]

  % =============== (a) campo uniforme ===============
  \begin{scope}
    \foreach \y in {-1.2,-0.6,0,0.6,1.2} {\draw[figvec] (-1.1,\y) -- (1.1,\y);}
    \node[figlbl, align=center] at (0,-2.1) {(a) uniforme:\\ paralelas y equiespaciadas};
  \end{scope}

  % =============== (b) carga puntual positiva ===============
  \begin{scope}[xshift=4.4cm]
    \foreach \a in {0,30,...,330} {\draw[figvec] (\a:0.3) -- (\a:1.45);}
    \node[figpos, minimum size=8pt] at (0,0) {};
    \node[figlbl, align=center] at (0,-2.1) {(b) carga $+$:\\ radiales, salientes};
  \end{scope}

  % =============== (c) dipolo ===============
  \begin{scope}[xshift=9.8cm]
    \coordinate (p) at (-0.85,0);
    \coordinate (n) at (0.85,0);
    % línea recta que las une
    \draw[figvec] (-0.55,0) -- (0.5,0);
    % arcos simétricos arriba y abajo
    \foreach \b in {25,50,75} {
      \draw[figvec] ($(p)+(\b:0.32)$) to[bend left=\b] ($(n)+(180-\b:0.32)$);
      \draw[figvec] ($(p)+(-\b:0.32)$) to[bend right=\b] ($(n)+(\b-180:0.32)$);
    }
    \node[figpos, minimum size=8pt] at (p) {};
    \node[figlbl, below=11pt] at (p) {$+Q$};
    \node[figneg, minimum size=8pt] at (n) {};
    \node[figlbl, below=11pt] at (n) {$-Q$};
    \node[figlbl, align=center] at (0,-2.1) {(c) dipolo:\\ nacen en $+Q$, mueren en $-Q$};
  \end{scope}

\end{tikzpicture}\\[0.5em]
{\small\itshape $\vec E$ es \textbf{tangente} a la línea en cada punto, y el
módulo es proporcional a la \textbf{densidad de líneas} por unidad de área
transversal. De ahí sale Coulomb sin cuentas: si $N$ líneas atraviesan la esfera
de radio $r$, entonces $|\vec E| \propto N/4\pi r^{2} \propto 1/r^{2}$, y el
$4\pi$ aparece solo, por geometría.}
\end{center}
